\section{Human Centric}\label{sec:human-centric}

A key difference between Industry 4.0 and Industry 5.0 is the inclusion of the human as a central part of the manufacturing process.
While Industry 4.0 focuses on mass production, mass productivity, and performance through automation, the goal is to eliminate or minimize the need for human intervention.
In Industry 5.0, a human-centric approach is the target, aiming to achieve the most effective collaboration between machines/robots and humans.
Robots will be tasked with repetitive and labor-intensive tasks, while humans are critical thinkers and creative cite{https://www.sciencedirect.com/science/article/pii/S2452414X21000558}.
Valuable in complex tasks, which are not easily automated.
cite{Industry 5.0: A survey on enabling technologies and potential applications}. 


Paper der skal læses:
- Industry 5.0: A survey on enabling technologies and potential applications
- A Literature Review of the Challenges and Opportunities of the Transition from Industry 4.0 to Society 5.0 
- Industry 5.0 in Manufacturing: Enhancing Resilience and Responsibility through AI-Driven Predictive Maintenance, Quality Control, and Supply Chain Optimization
- Digital Twins for Industry 5.0: Unlocking the Human Potential (Genlæs)
- A Digital Twin Driven Human-Centric Ecosystem for Industry 5.0 (Genlæs)
- Intelligent Manufacturing From the Perspective of Industry 5.0: Application Review and Prospects

\subsection{Enabling technologies}
%- Co-bots: Human robot collaboration
%- Extended Reality (XR): Virtual Reality(VR), Augmented Reality (AR), and Mixed Reality (MR). Help the operator in monitoring of the production line, new employee/operator training, assembly and maintenance instructions, and support. 
%- Digital Twins: Enable Decision Support, and Human/Machine Digital Twins collaboration.
